\title[Israel 7: Cristo--templo]{\bfseries{Israel, el pueblo de Dios}}
\subtitle{Estudio 7. De Cristo a la destrucción del templo}
\author{}\institute{}\date{}
\titlegraphic{\includegraphics[width=0.85\linewidth,height=0.38\textheight,keepaspectratio]{portada.jpg}\\[-1pt]
  {\fontsize{5}{6}\selectfont\color{sepia}\textit{Nota.} Elaboración propia: Sal 118:26 en texto hebreo consonántico (BHS).}}
\begin{document}
\frame{\titlepage}

% ===================================================
\section{Introducción}
\begin{frame}{¿Qué significó el año 70?}
\begin{columns}[T,onlytextwidth]
\begin{column}{0.47\textwidth}
\vspace{4pt}
\begin{cita}No quedará aquí piedra sobre piedra que no sea derribada.\fuente{(Mt 24:2)}\end{cita}
\vspace{4pt}
\begin{cita}¿Ha desechado Dios a su pueblo? ¡De ninguna manera! [\ldots] No ha desechado Dios a su pueblo, al cual conoció de antemano.\fuente{(Ro 11:1--2)}\end{cita}
\end{column}
\begin{column}{0.49\textwidth}
\centering
\begin{tikzpicture}[font=\scriptsize]
\node[bloqueoro,text width=2.6cm,minimum height=1.0cm] (a) at (-1.55,0) {«Dios rechazó\\a Israel»};
\node[bloqueoro,text width=2.6cm,minimum height=1.0cm] (b) at (1.55,0) {«Jesús nada tiene\\que ver con Israel»};
\begin{scope}[visible on=<2->]\foreach \n in {a,b}{\node[circle,fill=sepia,text=white,font=\tiny\bfseries,inner sep=1pt] at (\n.north east) {no};}\end{scope}
\begin{scope}[visible on=<3->]\node[etiq,text=sepia] at (0,-0.85) {\itshape dos extremos que los textos no avalan};\end{scope}
\begin{scope}[visible on=<4->]\node[bloqueazul,text width=5.7cm] (c) at (0,-1.55) {Lo que dicen los textos};\end{scope}
\begin{scope}[visible on=<5->]\node[bloque,text width=2.6cm,minimum height=1.5cm] (m) at (-1.55,-2.9) {\textbf{El Mesías}\\Jesús asume el llamado de Israel; un remanente judío es el núcleo};\end{scope}
\begin{scope}[visible on=<6->]\node[bloque,text width=2.6cm,minimum height=1.5cm] (j) at (1.55,-2.9) {\textbf{El juicio}\\ocurre dentro del pacto y lleva escrito un «hasta»};\end{scope}
\begin{scope}[visible on=<6->]\draw[flecha] (c.south -| m.north) -- (m.north);\end{scope}
\begin{scope}[visible on=<6->]\draw[flecha] (c.south -| j.north) -- (j.north);\end{scope}
\end{tikzpicture}
\elabpropia
\end{column}
\end{columns}
\end{frame}

\begin{frame}{Los ejes de la serie en esta etapa}
\centering
\begin{tikzpicture}
 \foreach \x/\n/\t/\d/\r [count=\ovc,evaluate=\ovc as \ovk using int(\ovc+0)] in {0/1/{Nación santa}/{¿Cómo encarna Jesús la vocación de Israel y quién la lleva adelante después de la Pascua?}/{Is 49:5--6; Mt 2:15; Ro 15:7--16; 1~P 2:9},
   4.3/2/{Restauración}/{¿Qué promete Jesús a Israel mientras anuncia la ruina del templo?}/{Mt 19:28; 23:39; Lc 21:24; 22:20; Hch 3:21},
   8.6/3/{Proyección}/{¿Significa el año 70 que Dios desechó a su pueblo?}/{Lv 26:44; Jer 31:35--37; Ro 11:1--2, 11, 29}}{\begin{scope}[visible on=<\ovk->]
  \node[caja,anchor=north] (b\n) at (\x,0) {\parbox[t][3.1cm][t]{3.55cm}{\RaggedRight
    {\small\textbf{\textcolor{ladrillo}{\t}}}\\[5pt]\small \d\par\vfill{\scriptsize\color{sepia}(\r)}}};
  \node[circle,fill=cafe,text=piedra,font=\scriptsize\bfseries,inner sep=1.6pt] at (b\n.north west) {\n};\end{scope}}
 \begin{scope}[visible on=<3->]\draw[flecha] (b1.east) -- (b2.west);\end{scope}
 \begin{scope}[visible on=<3->]\draw[flecha] (b2.east) -- (b3.west);\end{scope}
\end{tikzpicture}
\elabpropia
\vspace{6pt}
\begin{notalam}[Sobre las citas]
Las citas bíblicas son traducción propia desde el texto hebreo (BHS) y griego (NA28). Numeración española [hebrea] cuando difiere; en Jer 31, la de las ediciones judías tradicionales. En esta etapa el reino del norte ya no existe: «Israel» designa casi siempre al pueblo entero, y se indica cuando un texto evoca a las tribus del norte.
\end{notalam}
\end{frame}

% ===================================================
\section{Cronología}
\begin{frame}{Secuencia de la etapa}
\centering
\begin{tikzpicture}[x=0.78cm]
\draw[line width=1.6pt,cafe,-{Latex[length=3mm]}] (0,0) -- (15.7,0);
\foreach \x/\t/\b [count=\ovc,evaluate=\ovc as \ovk using int(\ovc+0)] in {0.6/{Nacimiento de Jesús}/{c. 6--4 a.\,e.\,c.},
                      2.65/{Juan y Jesús}/{c. 28--30},
                      4.7/{Cruz y Pentecostés}/{30 o 33},
                      6.75/{Concilio de Jerusalén}/{c. 48--49},
                      8.8/{Romanos}/{c. 57},
                      10.85/{Muerte de Jacobo}/{62},
                      12.9/{Revuelta}/{66},
                      14.95/{Templo; Masada}/{70; 73/74}}{\begin{scope}[visible on=<\ovk->]
  \fill[oro] (\x,0) circle (4pt);
  \node[above,font=\scriptsize\bfseries,text=ladrillo,text width=1.5cm,align=center] at (\x,0.2) {\t};
  \node[below,font=\tiny,text=sepia,text width=1.5cm,align=center] at (\x,-0.2) {\b};\end{scope}}
\begin{scope}[visible on=<8->]\draw[decorate,decoration={brace,amplitude=5pt,mirror},sepia] (0.4,-0.9) -- (4.9,-0.9) node[midway,below=5pt,font=\scriptsize\itshape]{El Mesías en Israel};\end{scope}
\begin{scope}[visible on=<8->]\draw[decorate,decoration={brace,amplitude=5pt,mirror},sepia] (5.6,-0.9) -- (11.1,-0.9) node[midway,below=5pt,font=\scriptsize\itshape]{Remanente y misión};\end{scope}
\begin{scope}[visible on=<8->]\draw[decorate,decoration={brace,amplitude=5pt,mirror},sepia] (12.6,-0.9) -- (15.3,-0.9) node[midway,below=5pt,font=\scriptsize\itshape]{Guerra judía};\end{scope}
\end{tikzpicture}
\notafig{Elaboración propia. Fechas aproximadas e.\,c., salvo indicación.}
\vspace{10pt}
\begin{columns}[T,onlytextwidth]
\begin{column}{0.48\textwidth}
\begin{notalam}[Fechas en debate]
La crucifixión se fecha el 7 de abril del 30 o el 3 de abril del 33. Josefo no fecha la caída de Masada; varios historiadores la sitúan en el 74 y no en el 73 (Mason, 2016).
\end{notalam}
\end{column}
\begin{column}{0.48\textwidth}
\begin{notalam}[Los evangelios y el 70]
La fecha de composición de Mateo y Lucas, antes o después del 70, se discute. Por eso conviene leer sus anuncios sobre el templo sin dar por resuelto ese punto.
\end{notalam}
\end{column}
\end{columns}
\end{frame}

% ===================================================
\section{Historia}
\begin{frame}{Judea bajo Roma}
\begin{columns}[T,onlytextwidth]
\begin{column}{0.52\textwidth}
\centering
\begin{tikzpicture}[font=\scriptsize]
\node[bloqueazul,text width=5.6cm] (r) at (0,0) {Roma: Pompeyo toma Jerusalén (63 a.\,e.\,c.)};
\node[bloque,text width=5.6cm] (p) at (0,-1.05) {Desde el 6 e.\,c., prefectos romanos en Cesarea};
\node[bloque,text width=5.6cm] (s) at (0,-2.2) {Sumo sacerdote (nombrado por Roma o los herodianos) y sanedrín};
\draw[flecha] (r) -- (p);
\draw[flecha] (p) -- (s);
\foreach \x/\t [count=\ovc,evaluate=\ovc as \ovk using int(\ovc+0)] in {-2.16/Fariseos,-0.72/Saduceos,0.72/Esenios,2.16/{«Cuarta\\filosofía»}}{\begin{scope}[visible on=<\ovk->]
  \node[bloqueoro,text width=1.1cm,minimum height=0.85cm,inner sep=2pt,font=\tiny] (g) at (\x,-3.6) {\t};
  \draw[oro!90!ladrillo,line width=0.6pt] (s.south) -- (g.north);\end{scope}}
\begin{scope}[visible on=<4->]\node[etiq,text=sepia] at (0,-4.3) {\itshape grupos del judaísmo (\emph{Ant.} 18.23)};\end{scope}
\end{tikzpicture}
\elabpropia
\end{column}
\begin{column}{0.44\textwidth}\footnotesize
\destacar{El templo de Herodes}\par
Rey vasallo de Roma, Herodes reconstruyó el templo; la obra seguía en curso en tiempos de Jesús (Jn 2:20).\par\medskip
\destacar{La «cuarta filosofía»}\par
Así llama Josefo a la corriente de la resistencia armada (\emph{Ant.} 18.23).\par\medskip
\begin{clave}
\textbf{La esperanza del pueblo}: «la consolación de Israel» (Lc 2:25) y «la redención de Jerusalén» (Lc 2:38).
\end{clave}
\end{column}
\end{columns}
\end{frame}

\begin{frame}{Josefo, testigo y apologista}
\begin{columns}[T,onlytextwidth]
\begin{column}{0.46\textwidth}
\centering
\begin{tikzpicture}[font=\scriptsize]
\foreach \i/\t in {0/{Sacerdote en Jerusalén},1/{Comandante judío en Galilea},2/{Se rinde en Jotapata (67)},3/{Anuncia a Vespasiano que será emperador},4/{Escribe en Roma bajo los Flavios}}{
  \node[bloque,text width=5.0cm,minimum height=0.55cm] (j\i) at (0,-0.95*\i) {\t};}
\foreach \i [evaluate=\i as \j using int(\i+1)] in {0,1,2,3}{\draw[flecha] (j\i) -- (j\j);}
\end{tikzpicture}
\notafig{Elaboración propia a partir de Josefo, \emph{Guerra} 3.400--402, y Mason (2016).}
\end{column}
\begin{column}{0.50\textwidth}\footnotesize
\destacar{Lo que registra de esta etapa}
\begin{itemize}
\item Juan el Bautista (\emph{Ant.} 18.116--119).
\item «Jacobo, hermano de Jesús, llamado Cristo», lapidado en el 62 (\emph{Ant.} 20.200).
\item La guerra desde dentro: facciones, asedio e incendio del templo.
\end{itemize}
\vspace{4pt}
\begin{notalam}[Leer con cautela]
Josefo es testigo directo de la guerra, pero también apologista de Tito y de la casa flavia. Sus cifras y los discursos que pone en boca de los protagonistas deben leerse con cautela (Mason, 2016).
\end{notalam}
\end{column}
\end{columns}
\end{frame}

\begin{frame}{La guerra judía (66--73/74)}
\centering
\begin{tikzpicture}
\draw[line width=1.4pt,cafe,-{Latex[length=2.6mm]}] (-1.6,0.55) -- (11.3,0.55);
 \foreach \x/\n/\a/\t/\d/\r [count=\ovc,evaluate=\ovc as \ovk using int(\ovc+0)] in {0/1/66/{Estallido}/{Floro saquea el tesoro; cesan los sacrificios por el emperador; Cestio Galo derrotado en Bet Horón.}/{\emph{Guerra} 2.293, 409},
   3.25/2/67--68/{Galilea}/{Vespasiano y Tito toman Jotapata y Gamla; Josefo se rinde.}/{\emph{Guerra} 3.400--402},
   6.5/3/70/{Asedio}/{Facciones rivales; grano quemado; hambre. El templo arde en Av (Josefo: día 10).}/{\emph{Guerra} 5.24--26; 6.250},
   9.75/4/71--74/{Triunfo y Masada}/{Despojos del templo en Roma (71). Masada cae en el 73 o 74.}/{\emph{Guerra} 7.148--150; 7.389--401}}{\begin{scope}[visible on=<\ovk->]
  \node[hito] at (\x,0.55) {};
  \node[font=\scriptsize\bfseries,text=ladrillo] at (\x,0.9) {\a};
  \node[bloque,anchor=north,text width=2.6cm,minimum height=2.45cm,align=left] (b\n) at (\x,0.2) {{\footnotesize\bfseries\color{ladrillo}\t}\\[2pt]\scriptsize \d\\[2pt]{\tiny\color{sepia}(\r)}};\end{scope}}
 \begin{scope}[visible on=<4->]\foreach \i/\j in {1/2,2/3,3/4}{\draw[flecha] (b\i.east) -- (b\j.west);}\end{scope}
\end{tikzpicture}
\elabpropia
\vspace{4pt}
\begin{columns}[T,onlytextwidth]
\begin{column}{0.48\textwidth}
\begin{notalam}[\emph{Fiscus Iudaicus}]
Tras el 70, las dos dracmas que cada judío pagaba al templo pasan a Júpiter Capitolino (\emph{Guerra} 7.218; Dion Casio 66.7.2).
\end{notalam}
\end{column}
\begin{column}{0.48\textwidth}
\begin{notalam}[¿Quiso Tito salvar el templo?]
Sí, según Josefo (\emph{Guerra} 6.236--243); no, según Sulpicio Severo (\emph{Crónica} 2.30). Las fuentes se contradicen.
\end{notalam}
\end{column}
\end{columns}
\end{frame}

\begin{frame}{El escenario de la guerra}
\begin{columns}[T,onlytextwidth]
\begin{column}{0.40\textwidth}\centering\includegraphics[width=\linewidth,height=0.76\textheight,keepaspectratio]{mapa_galilea.pdf}\par\vspace{1pt}{\fontsize{4.5}{5.5}\selectfont\color{sepia}\textit{Nota.} Elaboración propia (GSHHS/WDB, Basemap); ubicaciones aproximadas.\par}\end{column}
\begin{column}{0.56\textwidth}\scriptsize\RaggedRight
{\scriptsize\RaggedRight
\tikz[baseline=-0.6ex]\draw[ladrillo,line width=1.5pt,line cap=round](0,0)--(0.6,0);~Vespasiano, 67: Ptolemaida $\to$ Jotapata\par
\tikz[baseline=-0.6ex]\draw[cafe,line width=1.5pt,line cap=round,dash pattern=on 3.4pt off 1.8pt](0,0)--(0.6,0);~Tito, 70: Cesarea $\to$ Jerusalén\par
\tikz[baseline=-0.6ex]\draw[oro,line width=1.5pt,line cap=round,dash pattern=on 1.2pt off 1.4pt](0,0)--(0.6,0);~Huida a Pela (tradición; Eusebio, HE 3.5.3)\par}
\medskip

\textbf{\color{ladrillo}67}: Vespasiano somete Galilea; Jotapata y Gamla.\par\smallskip
\textbf{\color{ladrillo}70}: Tito sube de Cesarea y sitia Jerusalén; el templo arde el 9/10 de Av.\par\smallskip
\textbf{\color{ladrillo}73/74}: cae Masada.\par\smallskip
\textbf{\color{ladrillo}Huida a Pela} (Eusebio, HE 3.5.3): tradición discutida; cf. Lc 21:21.\par
\begin{notalam}[La cifra de Josefo]
\textbf{1.100.000 muertos} (\emph{Guerra} 6.420): excede cualquier estimación razonable de la población. Conviene citarla siempre «según Josefo».
\end{notalam}
\end{column}
\end{columns}
\end{frame}

\begin{frame}{Precisiones sobre el año 70}
\footnotesize El 70 suele resumirse así: más de un millón de muertos, fin del sacerdocio y del sanedrín, comienzo de la Diáspora.\par\medskip
\centering\scriptsize\renewcommand{\arraystretch}{1.3}
\begin{tabular}{>{\RaggedRight\arraybackslash}p{3.3cm}>{\RaggedRight\arraybackslash}p{8.6cm}}
\toprule
{\color{ladrillo}\bfseries Afirmación} & {\color{ladrillo}\bfseries Precisión}\\\midrule
Más de un millón de muertos & Cifra de Josefo (\emph{Guerra} 6.420), considerada exagerada.\\
Comienza la Diáspora & Ya existía desde el exilio babilónico (Hch 2:9--11). La expulsión de Jerusalén llega tras Bar Kojba (132--135): \emph{Aelia Capitolina}.\\
Fin del sacerdocio y del sanedrín & Cesan el sacrificio y el sumo sacerdocio; la tradición judía reconstituye un tribunal de sabios en Yavne.\\
Persecuciones posteriores & Expulsiones: Inglaterra 1290; Francia 1306 y 1394; España 1492. Peste negra: 1347--1351 (s. \textsc{xiv}).\\
\bottomrule
\end{tabular}
\elabpropia
\vspace{6pt}
\begin{clave}
Con estos matices, la idea central sigue en pie: el juicio fue grave, pero «suspendido, mas no para siempre» (Jer 31:35--37; Ro 11:25--26).
\end{clave}
\end{frame}

\begin{frame}{Arqueología: el mundo de Jesús}
\centering
\begin{tikzpicture}
 \foreach \x/\n/\t/\d/\r in {0/1/{Pilato}/{Cesarea, 1961: «[Po]ntius Pilatus [prae]fectus Iuda[eae]». Nombre y título del prefecto.}/{Inscripción latina},
   3.25/2/{Caifás}/{Osario con «Yehosef bar Qayafa» (1990). Identificación probable, lectura en debate.}/{\emph{Ant.} 18.35; Evans, 2003},
   6.5/3/{Siloé y Betesda}/{Gradas de Siloé (2004; Jn 9:7) y dos estanques junto a Santa Ana (Jn 5:2).}/{Reich, 2011},
   9.75/4/{Advertencia}/{Losa griega: ningún extranjero pase la balaustrada del santuario.}/{\emph{Guerra} 5.193; Ef 2:14}}{
  \node[caja,anchor=north] (b\n) at (\x,0) {\parbox[t][2.75cm][t]{2.65cm}{\RaggedRight
    {\small\textbf{\textcolor{ladrillo}{\t}}}\\[4pt]\scriptsize \d\par\vfill{\tiny\color{sepia}(\r)}}};
  \node[circle,fill=cafe,text=piedra,font=\scriptsize\bfseries,inner sep=1.6pt] at (b\n.north west) {\n};}
\end{tikzpicture}
\elabpropia
\vspace{4pt}
\begin{notalam}[Inscripción de Teódoto (anterior al 70)]
Una sinagoga en Jerusalén «para la lectura de la ley y la enseñanza de los mandamientos», con hospedaje para los que venían del extranjero (cf. Hch 6:9). La sinagoga ya convivía con el templo antes de su caída.
\end{notalam}
\end{frame}

\begin{frame}{Arqueología: la caída}
\centering
\begin{tikzpicture}
 \foreach \x/\n/\t/\d/\r in {0/1/{Calle herodiana}/{Sillares derribados en el 70 sobre el pavimento junto al arco de Robinson; piedra del «toque de trompeta».}/{Mt 24:2; \emph{Guerra} 4.582},
   3.25/2/{Casa Quemada}/{Vivienda incendiada en el 70; pesa con el nombre sacerdotal «Bar Kathros».}/{Avigad, 1983; Pesajim 57a},
   6.5/3/{Monedas}/{Rebeldes: «Jerusalén la santa», «por la redención de Sion». Roma: \textsc{Iudaea capta}.}/{Meshorer, 2001},
   9.75/4/{Arco de Tito}/{Roma, después del 81: candelabro, mesa y trompetas en la procesión triunfal.}/{\emph{Guerra} 7.148--150; Fine, 2016}}{
  \node[caja,anchor=north] (b\n) at (\x,0) {\parbox[t][2.75cm][t]{2.65cm}{\RaggedRight
    {\small\textbf{\textcolor{ladrillo}{\t}}}\\[4pt]\scriptsize \d\par\vfill{\tiny\color{sepia}(\r)}}};
  \node[circle,fill=cafe,text=piedra,font=\scriptsize\bfseries,inner sep=1.6pt] at (b\n.north west) {\n};}
\end{tikzpicture}
\elabpropia
\vspace{4pt}
\begin{notalam}[Masada: hallazgo y debate]
Yadin (1963--1965) excavó la rampa romana, la sinagoga, baños rituales y un rollo de Ez 35--38, la visión de los huesos secos, bajo el piso de la sinagoga. El suicidio colectivo solo lo narra Josefo; su exactitud se discute (Ben-Yehuda, 1995).
\end{notalam}
\end{frame}

% ===================================================
\section{El Mesías}
\begin{frame}{«La plenitud del tiempo»}
\begin{columns}[T,onlytextwidth]
\begin{column}{0.48\textwidth}
\begin{cita}Cuando vino la plenitud del tiempo, Dios envió a su Hijo, nacido de mujer, nacido bajo la ley.\fuente{(Gá 4:4)}\end{cita}
\vspace{2pt}
\footnotesize \destacar{Zacarías} (Lc 1:68--79): Dios «ha visitado y redimido a su pueblo [\ldots] para que le sirvamos (\grc{λατρεύειν}) en santidad».\par\smallskip
Redimido, Israel recupera su oficio sacerdotal.
\end{column}
\begin{column}{0.48\textwidth}
\centering
\begin{tikzpicture}[font=\scriptsize]
\node[bloqueazul,text width=2.6cm,font=\small] (p) at (0,0) {\grc{πλήρωμα}\\[-1pt]{\scriptsize\mdseries «plenitud»}};
\begin{scope}[visible on=<2->]\node[bloque,text width=2.3cm] (g) at (0,1.55) {\textbf{Gá 4:4}\\la plenitud del tiempo};\end{scope}
\begin{scope}[visible on=<3->]\node[bloque,text width=2.3cm] (i) at (-1.65,-1.6) {\textbf{Ro 11:12}\\la plenitud de Israel};\end{scope}
\begin{scope}[visible on=<4->]\node[bloque,text width=2.3cm] (n) at (1.65,-1.6) {\textbf{Ro 11:25}\\la plenitud de los gentiles};\end{scope}
\begin{scope}[visible on=<4->]\draw[flecha] (g) -- (p);\end{scope}
\begin{scope}[visible on=<4->]\draw[flecha] (p) -- (i);\end{scope}
\begin{scope}[visible on=<4->]\draw[flecha] (p) -- (n);\end{scope}
\end{tikzpicture}
\elabpropia
\vspace{2pt}
\begin{clave}La misma palabra une la venida del Hijo con lo que Pablo espera para Israel y las naciones.\end{clave}
\end{column}
\end{columns}
\end{frame}

\begin{frame}{Una voz en el desierto}
\begin{columns}[T,onlytextwidth]
\begin{column}{0.48\textwidth}
\centering{\large\heb{קוֹל קוֹרֵא בַּמִּדְבָּר פַּנּוּ דֶּרֶךְ יְהוָה}}\par\smallskip
\begin{cita}Una voz clama: En el desierto preparad el camino de YHWH.\fuente{(Is 40:3)}\end{cita}
\footnotesize Juan el Bautista: nuevo éxodo, consuelo (Is 40:1) y llamado a la \emph{teshuvá}.
\end{column}
\begin{column}{0.48\textwidth}
\centering
\begin{tikzpicture}[font=\scriptsize]
\node[bloqueazul,text width=1.0cm,minimum height=0.95cm] (tm) at (0,0) {TM};
\begin{scope}[visible on=<2->]\node[bloque,text width=1.35cm,minimum height=0.95cm,right=0.45cm of tm] (t1) {Una voz clama:};\end{scope}
\begin{scope}[visible on=<3->]\node[bloqueoro,text width=1.6cm,minimum height=0.95cm,right=0.2cm of t1] (t2) {en el desierto preparad};\end{scope}
\begin{scope}[visible on=<4->]\node[bloqueazul,text width=1.0cm,minimum height=0.95cm] (lx) at (0,-1.35) {LXX;\\Mc 1:3};\end{scope}
\begin{scope}[visible on=<5->]\node[bloqueoro,text width=1.6cm,minimum height=0.95cm,right=0.45cm of lx] (l1) {voz del que clama en el desierto:};\end{scope}
\begin{scope}[visible on=<6->]\node[bloque,text width=1.35cm,minimum height=0.95cm,right=0.2cm of l1] (l2) {preparad};\end{scope}
\begin{scope}[visible on=<6->]\draw[flecha] (tm) -- (t1);\end{scope} \begin{scope}[visible on=<6->]\draw[flecha] (lx) -- (l1);\end{scope}
\end{tikzpicture}
\notafig{Elaboración propia. En oro, dónde se ubica «en el desierto».}
\vspace{4pt}
\begin{notalam}[Precisión filológica]
El TM une «en el desierto» con «preparad»; la LXX y Mc 1:3, con «voz». Qumrán aplicaba el texto a su propio retiro en el desierto (1QS 8.12--16).
\end{notalam}
\end{column}
\end{columns}
\end{frame}

\begin{frame}{Simeón y Ana}
\fotolado{panel_is49.jpg}{Elaboración propia: Is 49:6b en texto hebreo consonántico (BHS).}
\begin{columns}[c,onlytextwidth]
\begin{column}{0.57\textwidth}
\begin{cita}Poco es que me seas siervo para levantar las tribus de Jacob [\ldots]; te daré por luz de las naciones, para que seas mi salvación hasta el extremo de la tierra.\fuente{(Is 49:6)}\end{cita}
\begin{cita}Luz para revelación a las naciones y gloria de tu pueblo Israel.\fuente{(Lc 2:32)}\end{cita}
\scriptsize
\begin{itemize}
\item Las dos misiones del siervo, juntas: \destacar{gloria para Israel y luz para las naciones}.
\item «Para caída y levantamiento de muchos en Israel» (2:34; cf. Ro 11:11).
\item Ana, «de la tribu de Aser»: el antiguo reino del norte presente en el templo (2:36; cf. 2~Cr 30:11).
\end{itemize}
\end{column}
\begin{column}{0.41\textwidth}\end{column}
\end{columns}
\end{frame}

\begin{frame}{Nazaret: el año del jubileo}
\begin{columns}[T,onlytextwidth]
\begin{column}{0.48\textwidth}
\centering{\heb{לִקְרֹא לִשְׁבוּיִם דְּרוֹר}}\par\smallskip
\begin{cita}Me ha enviado a proclamar libertad a los cautivos [\ldots], a proclamar el año agradable del Señor. Hoy se ha cumplido esta Escritura.\fuente{(Lc 4:18--21; Is 61:1--2)}\end{cita}
\end{column}
\begin{column}{0.48\textwidth}
\centering
\begin{tikzpicture}[font=\scriptsize]
\node[bloque,text width=1.65cm,minimum height=1.25cm] (a) at (0,0) {proclamar\\libertad\\(\emph{derór})};
\begin{scope}[visible on=<2->]\node[bloque,text width=1.65cm,minimum height=1.25cm] (b) at (2.05,0) {el año\\agradable\\del Señor};\end{scope}
\begin{scope}[visible on=<3->]\node[draw=sepia,dashed,rounded corners=3pt,text=sepia,text width=1.65cm,minimum height=1.25cm,align=center,inner sep=4pt] (c) at (4.3,0) {el día de\\venganza\\(Is 61:2b)};\end{scope}
\begin{scope}[visible on=<3->]\draw[flecha] (a) -- (b);\end{scope}
\begin{scope}[visible on=<3->]\draw[cafeosc,line width=1.6pt] (3.17,-0.85) -- (3.17,0.85);\end{scope}
\begin{scope}[visible on=<4->]\node[etiq,text=cafeosc,font=\tiny\bfseries] at (3.17,1.05) {Jesús se detiene};\end{scope}
\begin{scope}[visible on=<4->]\draw[decorate,decoration={brace,amplitude=4pt,mirror},sepia] (-0.85,-0.8) -- (2.9,-0.8) node[midway,below=5pt,etiq]{«Hoy se ha cumplido» (Lc 4:21)};\end{scope}
\end{tikzpicture}
\elabpropia
\vspace{4pt}
\begin{notalam}[\heb{דְּרוֹר}, \emph{derór}]
Es la libertad del jubileo (Lv 25:10): los cautivos vuelven y la tierra se restituye. Jesús lee Is 61 como anuncio de ese año de liberación.
\end{notalam}
\end{column}
\end{columns}
\end{frame}

\begin{frame}{Los Doce y las doce tribus}
\begin{columns}[T,onlytextwidth]
\begin{column}{0.48\textwidth}
\begin{cita}Os sentaréis sobre doce tronos, para juzgar a las doce tribus de Israel.\fuente{(Mt 19:28; Lc 22:30)}\end{cita}
\footnotesize
\begin{itemize}
\item Matías completa el número antes de Pentecostés (Hch 1:15--26).
\item «Nuestras doce tribus» (Hch 26:7); Santiago escribe «a las doce tribus que están en la dispersión» (Stg 1:1).
\item «Las ovejas perdidas de la casa de Israel» (Mt 10:6; 15:24; cf. Ez 34).
\end{itemize}
\end{column}
\begin{column}{0.48\textwidth}
\centering
\begin{tikzpicture}[font=\scriptsize]
\node[bloque,text width=2.2cm] (n) at (-1.4,0) {Tribus del norte\\dispersas (722)};
\begin{scope}[visible on=<2->]\node[bloque,text width=2.2cm] (s) at (1.4,0) {Judá\\en la tierra};\end{scope}
\begin{scope}[visible on=<3->]\node[bloqueazul,text width=5.0cm] (u) at (0,-1.35) {Una sola nación (Ez 37:15--28)};\end{scope}
\begin{scope}[visible on=<3->]\draw[flecha] (n) -- (u.north -| n);\end{scope} \begin{scope}[visible on=<3->]\draw[flecha] (s) -- (u.north -| s);\end{scope}
\begin{scope}[visible on=<4->]\foreach \k in {0,...,11}{\node[hito] at (-2.2+\k*0.4,-2.15) {};}\end{scope}
\begin{scope}[visible on=<5->]\node[etiq] at (0,-2.5) {los Doce: signo de \emph{todo} Israel reconstituido (Sanders, 1985)};\end{scope}
\end{tikzpicture}
\elabpropia
\vspace{4pt}
\begin{notalam}[«Casa de Israel» en Mateo]
Es el pueblo judío de Galilea y Judea, no el antiguo reino del norte: a los Doce se les prohíbe entrar en Samaria (Mt 10:5--6).
\end{notalam}
\end{column}
\end{columns}
\end{frame}

\begin{frame}{Jesús en las fiestas de Israel}
\centering
\begin{tikzpicture}
 \foreach \x/\n/\t/\d [count=\ovc,evaluate=\ovc as \ovk using int(\ovc+0)] in {0/1/{Sucot (Jn 7--8)}/{«Si alguno tiene sed, venga a mí» (7:37), en la libación del agua; «Yo soy la luz del mundo» (8:12), junto a los candelabros del atrio.},
   4.3/2/{Janucá (Jn 10:22)}/{\grc{τὰ ἐγκαίνια}, «la Dedicación» (\heb{חֲנֻכָּה}): «era invierno, y Jesús andaba en el templo, en el pórtico de Salomón».},
   8.6/3/{Pésaj y el Hallel}/{Sal 118 se canta en Pésaj y Sucot; con el \emph{lulav} se aclama «\emph{hoshía na}» (Sucá 3:9). La multitud lo grita en la entrada (Mt 21:9).}}{\begin{scope}[visible on=<\ovk->]
  \node[bloque,anchor=north,text width=3.55cm,minimum height=2.35cm,align=left] (f\n) at (\x,0) {{\footnotesize\bfseries\color{ladrillo}\t}\\[3pt]\scriptsize\d};\end{scope}}
 \begin{scope}[visible on=<3->]\draw[flecha] (f1.east) -- (f2.west);\end{scope}
 \begin{scope}[visible on=<3->]\draw[flecha] (f2.east) -- (f3.west);\end{scope}
\end{tikzpicture}
\elabpropia
\vspace{4pt}
\centering{\large\heb{בָּרוּךְ הַבָּא בְּשֵׁם יְהוָה}}\par
{\footnotesize\itshape\color{sepia}«Bendito el que viene en el nombre de YHWH» (Sal 118:26): un mismo salmo para la entrada, la parábola y el lamento.}\par
\vspace{4pt}
\begin{notalam}[Tradición judía]
La Mishná describe la libación del agua y los grandes candelabros encendidos en el atrio de las mujeres durante Sucot (Sucá 5:1--3). Janucá recuerda la purificación del templo por los macabeos (Estudio~6).
\end{notalam}
\end{frame}

\begin{frame}{Los labradores: Mt 21:43}
\begin{columns}[T,onlytextwidth]
\begin{column}{0.42\textwidth}
\begin{cita}El reino de Dios será quitado de vosotros y será dado a un pueblo que produzca sus frutos.\fuente{(Mt 21:43)}\end{cita}
\footnotesize Antes, en el templo: «casa de oración [\ldots] para todos los pueblos» (Is 56:7) y «cueva de ladrones» (Jer 7:11), eco del sermón que anunció la caída del Primer Templo.
\end{column}
\begin{column}{0.54\textwidth}
\centering
\begin{tikzpicture}[font=\scriptsize]
\node[bloqueazul,text width=5.6cm] (v) at (0,0) {La viña es Israel (Is 5:7): no se destruye};
\begin{scope}[visible on=<2->]\node[bloqueoro,text width=2.25cm,minimum height=1.45cm] (a) at (-1.75,-1.55) {\textbf{Labradores}\\sacerdotes y fariseos: «hablaba de ellos» (21:45)};\end{scope}
\begin{scope}[visible on=<3->]\node[bloque,text width=2.25cm,minimum height=1.45cm] (b) at (1.75,-1.55) {\textbf{Otros labradores} (21:41)\\\grc{ἔθνει}, «a un pueblo», singular; núcleo: los Doce};\end{scope}
\begin{scope}[visible on=<3->]\draw[flecha] (a) -- (b);\end{scope}
\begin{scope}[visible on=<3->]\draw[ladrillo,line width=0.6pt] (v.south -| a) -- (a);\end{scope} \begin{scope}[visible on=<3->]\draw[ladrillo,line width=0.6pt] (v.south -| b) -- (b);\end{scope}
\begin{scope}[visible on=<4->]\node[bloque,text width=5.6cm] (p) at (0,-3.5) {La piedra desechada (Sal 118:22) es el Hijo; Pedro la aplica a los «edificadores» (Hch 4:11)};\end{scope}
\end{tikzpicture}
\elabpropia
\vspace{2pt}
\begin{clave}\textbf{Cambio de liderazgo, no de pueblo} (Snodgrass, 2008). No dice \grc{τοῖς ἔθνεσιν}, «a los gentiles».\end{clave}
\end{column}
\end{columns}
\end{frame}

\begin{frame}{El lamento: Mt 23:39}
\fotolado{panel_mt23.jpg}{Elaboración propia: Mt 23:39 en escritura uncial con \emph{nomen sacrum} (texto NA28).}
\begin{columns}[c,onlytextwidth]
\begin{column}{0.57\textwidth}
\begin{cita}¡Jerusalén, Jerusalén [\ldots]! ¡Cuántas veces quise juntar a tus hijos [\ldots], y no quisisteis! He aquí, vuestra casa os es dejada desierta. Porque os digo que no me veréis desde ahora, \textbf{hasta que} digáis: Bendito el que viene en el nombre del Señor.\fuente{(Mt 23:37--39)}\end{cita}
\scriptsize
\begin{itemize}
\item «Vuestra casa»: el templo (Jer 12:7; 22:5).
\item \grc{ἕως ἂν εἴπητε}: condición y plazo; Jerusalén lo verá de nuevo cuando lo reciba (Allison, 1983).
\item Sal 118:26: el saludo de los peregrinos y de los sacerdotes «desde la casa de YHWH».
\end{itemize}
\end{column}
\begin{column}{0.41\textwidth}\end{column}
\end{columns}
\end{frame}

\begin{frame}{Lc 21:24: «hasta que»}
\fotolado{panel_lc21.jpg}{Elaboración propia: Lc 21:24b en escritura uncial (texto NA28).}
\begin{columns}[c,onlytextwidth]
\begin{column}{0.57\textwidth}
\begin{cita}Serán llevados cautivos a todas las naciones; y Jerusalén será hollada por los gentiles, hasta que los tiempos de los gentiles se cumplan.\fuente{(Lc 21:24)}\end{cita}
\scriptsize
\begin{itemize}
\item Lenguaje de Dt 28:64 y Jeremías; «hollada»: posible eco de Zac 12:3 LXX.
\item \grc{ἄχρι οὗ}: la misma construcción de Ro 11:25, «hasta que haya entrado la plenitud de los gentiles».
\item «Tiempos de los gentiles»: ¿dominio, misión o juicio? En debate; el límite, en cambio, es claro.
\end{itemize}
\precision{la huida a Pela (Eusebio, HE 3.5.3) coincide con Lc 21:21; es tradición discutida.}
\end{column}
\begin{column}{0.41\textwidth}\end{column}
\end{columns}
\end{frame}

\begin{frame}{La copa del nuevo pacto}
\fotolado{panel_jer31.jpg}{Elaboración propia: Jer 31:31 [31:30] en texto hebreo consonántico (BHS).}
\begin{columns}[c,onlytextwidth]
\begin{column}{0.57\textwidth}
\begin{cita}Haré un nuevo pacto con la casa de Israel y con la casa de Judá.\fuente{(Jer 31:31 [31:30])}\end{cita}
\begin{cita}Esta copa es el nuevo pacto en mi sangre, que por vosotros se derrama.\fuente{(Lc 22:20)}\end{cita}
\scriptsize
\begin{itemize}
\item \heb{בְּרִית חֲדָשָׁה} $\to$ \grc{διαθήκη καινή} (Jer 38:31 LXX) $\to$ Lc 22:20; 1~Co 11:25.
\item «Casa de Israel» (norte) y «casa de Judá» (sur): un solo pacto para los dos reinos.
\item «La sangre del pacto» (Mt 26:28; Éx 24:8). Heb 8:8--12 cita Jer 31 completo.
\end{itemize}
\end{column}
\begin{column}{0.41\textwidth}\end{column}
\end{columns}
\end{frame}

\begin{frame}{«Mirarán al que traspasaron»}
\centering
\begin{tikzpicture}[node distance=0.9cm]
\node[caja,text width=3.3cm,align=center] (h) {\textbf{TM (Zac 12:10)}\\[3pt]\heb{אֵת אֲשֶׁר־דָּקָרוּ}\\[2pt]\emph{daqáru}\\[2pt]«traspasaron»};
\begin{scope}[visible on=<2->]\node[caja,text width=3.3cm,align=center,right=of h] (g) {\textbf{LXX}\\[3pt]\grc{ἀνθ᾽ ὧν κατωρχήσαντο}\\[2pt]«se burlaron danzando»\\(leyó \heb{רקדו})};\end{scope}
\begin{scope}[visible on=<3->]\node[caja,text width=3.3cm,align=center,right=of g] (p) {\textbf{Juan 19:37}\\[3pt]\grc{ὄψονται εἰς ὃν ἐξεκέντησαν}\\[2pt]sigue al hebreo,\\como Teodoción};\end{scope}
\begin{scope}[visible on=<3->]\draw[-{Latex[length=2.5mm]},oro!90!ladrillo,line width=1.4pt] (h) -- (g);\end{scope}
\begin{scope}[visible on=<3->]\draw[-{Latex[length=2.5mm]},oro!90!ladrillo,line width=1.4pt] (g) -- (p);\end{scope}
\end{tikzpicture}
\elabpropia
\vspace{6pt}
\begin{columns}[T,onlytextwidth]
\begin{column}{0.52\textwidth}\footnotesize
Hacen duelo «la casa de David y los habitantes de Jerusalén» (Judá), antesala del «manantial abierto [\ldots] para la purificación» (Zac 13:1). La serie lo lee como el momento en que es quitado el velo de la casa de Judá.
\end{column}
\begin{column}{0.44\textwidth}
\begin{notalam}[Precisión textual]
El TM dice «a mí» y luego «por él»; algunos manuscritos hebreos leen «a él».
\end{notalam}
\end{column}
\end{columns}
\end{frame}

\begin{frame}{¿Restaurarás el reino a Israel?}
\begin{columns}[T,onlytextwidth]
\begin{column}{0.50\textwidth}
\begin{cita}Señor, ¿restaurarás el reino a Israel en este tiempo? Y les dijo: No os toca a vosotros saber los tiempos [\ldots]; pero me seréis testigos en Jerusalén, en toda Judea, en Samaria y hasta el extremo de la tierra.\fuente{(Hch 1:6--8)}\end{cita}
\begin{clave}Jesús deja en pie el \textbf{contenido} de la pregunta; corrige el \textbf{calendario}.\end{clave}
\begin{notalam}[\grc{ἀποκαθιστάνεις}]
Es el verbo de Elías en Mal 3:23 [4:6] LXX y Mt 17:11: «restaurar».
\end{notalam}
\end{column}
\begin{column}{0.46\textwidth}
\centering
\begin{tikzpicture}[font=\scriptsize]
\fill[piedra!60!white,draw=ladrillo!60] (0,0) circle (2.15);
\fill[oro!14!white,draw=oro!90!ladrillo] (0,0) circle (1.55);
\fill[piedra,draw=ladrillo] (0,0) circle (0.98);
\fill[ladrillo] (0,0) circle (0.48);
\node[text=white,font=\tiny\bfseries] at (0,0) {Jerusalén};
\begin{scope}[visible on=<2->]\node[font=\tiny\bfseries,text=ladrillo] at (0,0.72) {Judea};\end{scope}
\begin{scope}[visible on=<3->]\node[font=\tiny\bfseries,text=cafeosc,align=center] at (0,1.27) {Samaria};\end{scope}
\begin{scope}[visible on=<4->]\node[font=\tiny\bfseries,text=ladrillo,align=center] at (0,1.86) {extremo de la tierra};\end{scope}
\begin{scope}[visible on=<5->]\node[etiq,text=sepia,text width=4.5cm] at (0,-2.5) {Samaria: región del antiguo reino del norte (Hch 8). El mapa es el de Is 49:6: \grc{ἕως ἐσχάτου τῆς γῆς}.};\end{scope}
\end{tikzpicture}
\elabpropia
\end{column}
\end{columns}
\end{frame}

% ===================================================
\section{El remanente}
\begin{frame}{Pentecostés y los miles}
\begin{columns}[T,onlytextwidth]
\begin{column}{0.46\textwidth}\footnotesize
\destacar{Shavuot: la diáspora reunida}\par\smallskip
«Partos, medos, elamitas, habitantes de Mesopotamia\ldots» (Hch 2:9--11).\par\smallskip
Pedro cita Joel 2:28--32 [3:1--5]: «derramaré mi Espíritu sobre toda carne».\par
\begin{notalam}[El norte y el sur en Hch 2]
Medos y Mesopotamia: destino de los deportados del reino del norte (2~R 17:6); Babilonia: el de Judá. Para Lucas, Pentecostés inaugura la reunión del Israel disperso (Jervell, 1972).
\end{notalam}
\end{column}
\begin{column}{0.50\textwidth}
\centering
\begin{tikzpicture}[font=\scriptsize]
\foreach \i/\v/\h in {0/{c. 3.000 personas}/{Hch 2:41},1/{c. 5.000 hombres}/{Hch 4:4},2/{«Gran multitud de sacerdotes»}/{Hch 6:7},3/{«Decenas de miles» (\grc{μυριάδες})}/{Hch 21:20}}{
  \node[bloque,text width=4.3cm,inner sep=3pt] (m\i) at (0,-0.95*\i) {\textbf{\v}\\[-1pt]{\tiny\color{sepia}\h}};}
\draw[flecha] (-2.85,0.2) -- (-2.85,-3.05) node[midway,etiq,rotate=90,above=1pt]{crecimiento};
\end{tikzpicture}
\elabpropia
\vspace{2pt}
\begin{clave}«Todos son celosos por la ley» (Hch 21:20); oran en el templo (2:46; 3:1).\end{clave}
\end{column}
\end{columns}
\end{frame}

\begin{frame}{De todas las naciones}
\begin{columns}[T,onlytextwidth]
\begin{column}{0.62\textwidth}\centering\includegraphics[width=\linewidth,height=0.80\textheight,keepaspectratio]{mapa_pentecostes.pdf}\par\vspace{1pt}{\fontsize{4.5}{5.5}\selectfont\color{sepia}\textit{Nota.} Elaboración propia (GSHHS/WDB, Basemap); ubicaciones aproximadas.\par}\end{column}
\begin{column}{0.35\textwidth}\scriptsize\RaggedRight
{\scriptsize\RaggedRight
\definecolor{mapacarmesi}{HTML}{D62839}
\definecolor{mapanaranja}{HTML}{F28C28}
\definecolor{mapavioleta}{HTML}{7B2CBF}
\definecolor{mapaturquesa}{HTML}{0F8B8D}
\definecolor{mapaazul}{HTML}{1D4E89}
\definecolor{mapaocre}{HTML}{B5651D}
\makebox[0.62cm][c]{\tikz[baseline=-0.6ex]\node[circle,fill=white,draw=black!80,line width=0.8pt,inner sep=1.2pt]{};}~Procedencia de los oyentes (Hch 2:9–11)\par}
\medskip
\begin{cita}[fontupper=\scriptsize\itshape]Moraban en Jerusalén judíos, varones piadosos, de todas las naciones bajo el cielo.\fuente{(Hch 2:5)}\end{cita}
Partos, medos, elamitas; Mesopotamia, Judea, Capadocia, Ponto, Asia, Frigia, Panfilia; Egipto, Libia, Roma; cretenses y árabes (Hch 2:9--11).\par\smallskip
El primer anuncio alcanza a la diáspora congregada en Shavuot.
\end{column}
\end{columns}
\end{frame}

\begin{frame}{Jacobo y la tienda de David}
\begin{columns}[T,onlytextwidth]
\begin{column}{0.42\textwidth}
\centering
\begin{tikzpicture}[font=\scriptsize]
\foreach \i/\t in {0/{Testigo de la resurrección (1~Co 15:7)},1/{«Columna» de la iglesia (Gá 2:9)},2/{Preside el concilio (Hch 15)},3/{Lapidado en el 62 por orden de Anán (\emph{Ant.} 20.200)}}{
  \node[bloque,text width=4.2cm,align=left] (k\i) at (0,-0.98*\i) {\t};}
\foreach \i [evaluate=\i as \j using int(\i+1)] in {0,1,2}{\draw[flecha] (k\i) -- (k\j);}
\node[etiq,text=sepia] at (0,-3.55) {«el Justo» (Hegesipo, en Eusebio, HE 2.23)};
\end{tikzpicture}
\elabpropia
\end{column}
\begin{column}{0.54\textwidth}
\centering{\large\heb{אָקִים אֶת־סֻכַּת דָּוִיד הַנֹּפֶלֶת}}\par\smallskip
\begin{cita}Reedificaré la tienda de David, que está caída [\ldots], para que el resto de los hombres busque al Señor, y todas las naciones sobre las cuales es invocado mi nombre.\fuente{(Hch 15:16--17; Am 9:11--12 LXX)}\end{cita}
\begin{notalam}[TM y LXX en Am 9:12]
El TM dice «posean el resto de \emph{Edom}» (\heb{אֱדוֹם}); la LXX leyó \emph{adam}, «hombres», y «busquen». Los gentiles entran sin dejar de serlo.
\end{notalam}
\end{column}
\end{columns}
\end{frame}

\begin{frame}{«Al judío primeramente»}
\begin{columns}[T,onlytextwidth]
\begin{column}{0.46\textwidth}
\begin{cita}Poder de Dios para salvación a todo el que cree, al judío primeramente, y también al griego.\fuente{(Ro 1:16)}\end{cita}
\begin{cita}Por la esperanza de Israel estoy sujeto con esta cadena.\fuente{(Hch 28:20)}\end{cita}
\begin{notalam}[Dos lecturas de Hechos]
¿Historia trágica (Tannehill, 1985) o Israel dividido con un remanente numeroso (Jervell, 1972)? Hechos termina con la puerta abierta (28:30).
\end{notalam}
\end{column}
\begin{column}{0.50\textwidth}
\centering
\begin{tikzpicture}[font=\scriptsize]
\node[bloque,text width=4.6cm] (a) at (0,0) {Ro 1:16: «al judío primeramente, y también al griego»};
\begin{scope}[visible on=<2->]\node[bloqueoro,text width=4.6cm] (b) at (0,-1.15) {Antioquía: «nos volvemos a los gentiles», con Is 49:6 (Hch 13:46--47)};\end{scope}
\begin{scope}[visible on=<3->]\node[bloque,text width=4.6cm] (c) at (0,-2.3) {En la ciudad siguiente, otra vez a la sinagoga (14:1)};\end{scope}
\begin{scope}[visible on=<4->]\node[bloqueazul,text width=4.6cm] (d) at (0,-3.45) {Roma: unos creen y otros no (28:24); Pablo cita Is 6:9--10};\end{scope}
\begin{scope}[visible on=<4->]\draw[flecha] (a) -- (b);\end{scope} \begin{scope}[visible on=<4->]\draw[flecha] (b) -- (c);\end{scope} \begin{scope}[visible on=<4->]\draw[flecha] (c) -- (d);\end{scope}
\begin{scope}[visible on=<5->]\node[etiq,text=sepia] at (0,-4.25) {\grc{Ἰουδαίῳ τε πρῶτον καὶ Ἕλληνι}: el orden se repite};\end{scope}
\end{tikzpicture}
\elabpropia
\end{column}
\end{columns}
\end{frame}

\begin{frame}{Velo en parte}
\begin{columns}[T,onlytextwidth]
\begin{column}{0.46\textwidth}
\centering
\begin{tikzpicture}[font=\scriptsize]
\node[bloque,text width=4.4cm,align=left] (a) at (0,0) {\textbf{Hch 4:11}: la piedra desechada, hecha cabeza del ángulo (Sal 118:22)};
\begin{scope}[visible on=<2->]\node[bloque,text width=4.4cm,align=left] (b) at (0,-1.2) {\textbf{Mt 23:39}: «no me veréis [\ldots] hasta que digáis: Bendito el que viene»};\end{scope}
\begin{scope}[visible on=<3->]\node[bloque,text width=4.4cm,align=left] (c) at (0,-2.4) {\textbf{2~Co 3:14}: «hasta el día de hoy les queda el mismo velo»};\end{scope}
\begin{scope}[visible on=<4->]\node[bloqueazul,text width=4.4cm] (d) at (0,-3.6) {Un velo con límite: «cuando se vuelva al Señor»};\end{scope}
\begin{scope}[visible on=<4->]\draw[flecha] (c) -- (d);\end{scope}
\end{tikzpicture}
\elabpropia
\end{column}
\begin{column}{0.50\textwidth}
\begin{cita}Hasta hoy, cada vez que se lee a Moisés, el velo está puesto sobre su corazón. Pero cuando se vuelva al Señor, el velo es quitado.\fuente{(2~Co 3:15--16)}\end{cita}
\begin{notalam}[Griego y hebreo]
\grc{ἐπωρώθη} (2~Co 3:14) y \grc{πώρωσις ἀπὸ μέρους} (Ro 11:25) comparten raíz. \grc{ἡνίκα ἐὰν ἐπιστρέψῃ} viene de Éx 34:34 LXX; \grc{ἐπιστρέφω} traduce \heb{שׁוּב}, el verbo del regreso. La «lectura del antiguo pacto» es la lectura de Moisés en la sinagoga.
\end{notalam}
\end{column}
\end{columns}
\end{frame}

\begin{frame}{Razones del endurecimiento}
\centering
\begin{tikzpicture}
 \foreach \x/\n/\t/\d/\r [count=\ovc,evaluate=\ovc as \ovk using int(\ovc+0)] in {0/1/{Soberanía}/{El endurecimiento sirve al plan de salvación de los gentiles.}/{Ro 9:18; 11:11},
   3.25/2/{Ignorancia}/{«Ignorando la justicia de Dios» (\grc{ἀγνοοῦντες}); «el fin de la ley es Cristo».}/{Ro 10:3--4},
   6.5/3/{«Soberbia»}/{Establecer la propia justicia; pero Pablo reconoce «celo de Dios, no conforme a conocimiento».}/{Ro 10:2--3},
   9.75/4/{Temporalidad}/{No es definitivo: «en parte» y «hasta». «Todo Israel será salvo».}/{Ro 11:25--26}}{\begin{scope}[visible on=<\ovk->]
  \node[caja,anchor=north] (b\n) at (\x,0) {\parbox[t][2.8cm][t]{2.65cm}{\RaggedRight
    {\small\textbf{\textcolor{ladrillo}{\t}}}\\[4pt]\scriptsize \d\par\vfill{\tiny\color{sepia}(\r)}}};
  \node[circle,fill=cafe,text=piedra,font=\scriptsize\bfseries,inner sep=1.6pt] at (b\n.north west) {\n};\end{scope}}
 \begin{scope}[visible on=<4->]\draw[oro,line width=1.4pt,rounded corners=3pt] ($(b4.north west)+(-0.08,0.08)$) rectangle ($(b4.south east)+(0.08,-0.08)$);\end{scope}
 \begin{scope}[visible on=<4->]\node[circle,fill=cafe,text=piedra,font=\scriptsize\bfseries,inner sep=1.6pt] at (b4.north west) {4};\end{scope}
\end{tikzpicture}
\notafig{Elaboración propia. En oro, el aspecto que esta etapa pone de relieve.}
\vspace{6pt}
\begin{notalam}[Precisión]
La «ignorancia» está en Ro 10:3 (10:4 da la razón). Antes de hablar de soberbia, Pablo reconoce el celo de Israel (10:2).
\end{notalam}
\end{frame}

% ===================================================
\section{Vocación}
\begin{frame}{Jesús, el Israel fiel}
\begin{columns}[T,onlytextwidth]
\begin{column}{0.40\textwidth}
\centering{\heb{וּמִמִּצְרַיִם קָרָאתִי לִבְנִי}}\par\smallskip
\begin{cita}De Egipto llamé a mi hijo.\fuente{(Os 11:1; Mt 2:15)}\end{cita}
\begin{notalam}[Mateo sigue el hebreo]
Mateo lee «mi hijo», como el TM; la LXX dice «sus hijos». En Oseas el hijo es el pueblo (Éx 4:22).
\end{notalam}
\end{column}
\begin{column}{0.56\textwidth}
\centering
\begin{tikzpicture}[font=\scriptsize]
\node[bloqueazul,text width=2.7cm] at (-1.6,0.1) {Israel};
\node[bloqueazul,text width=2.7cm] at (1.6,0.1) {Jesús};
\foreach \y/\l/\r [count=\ovc,evaluate=\ovc as \ovk using int(\ovc+0)] in {-0.85/{Hijo llamado de Egipto (Éx 4:22; Os 11:1)}/{Huida y regreso de Egipto (Mt 2:15)},
                      -1.95/{Cuarenta años en el desierto}/{Cuarenta días en el desierto (Mt 4:1--11)},
                      -3.05/{Donde Israel murmuró y cayó}/{el Hijo confía: Dt 8:3; 6:16; 6:13}}{\begin{scope}[visible on=<\ovk->]
  \node[bloqueoro,text width=2.7cm,minimum height=0.85cm] (l) at (-1.6,\y) {\l};
  \node[bloque,text width=2.7cm,minimum height=0.85cm] (r) at (1.6,\y) {\r};
  \draw[flecha] (l) -- (r);\end{scope}}
\end{tikzpicture}
\elabpropia
\end{column}
\end{columns}
\end{frame}

\begin{frame}{El siervo de Isaías}
\begin{columns}[T,onlytextwidth]
\begin{column}{0.50\textwidth}
\centering
\begin{tikzpicture}[font=\scriptsize]
\node[bloque,text width=2.0cm,minimum height=1.2cm] (a) at (-1.55,0) {El siervo se llama\\«Israel»\\(Is 49:3)};
\begin{scope}[visible on=<2->]\node[bloque,text width=2.0cm,minimum height=1.2cm] (b) at (1.55,0) {Su misión: hacer volver a Jacob\\(Is 49:5)};\end{scope}
\begin{scope}[visible on=<2->]\draw[oro!90!ladrillo,line width=1pt,<->] (a) -- (b);\end{scope}
\begin{scope}[visible on=<3->]\node[etiq,text=sepia] at (0,0.95) {\itshape tensión};\end{scope}
\begin{scope}[visible on=<4->]\node[bloqueazul,text width=5.2cm] (c) at (0,-1.75) {El Mesías encarna a Israel para restaurarlo};\end{scope}
\begin{scope}[visible on=<4->]\draw[flecha] (a.south) -- (c.north -| a.south);\end{scope}
\begin{scope}[visible on=<4->]\draw[flecha] (b.south) -- (c.north -| b.south);\end{scope}
\begin{scope}[visible on=<5->]\node[bloqueoro,text width=5.2cm] (d) at (0,-3.05) {Éx 19:6 se cumple primero en una persona};\end{scope}
\begin{scope}[visible on=<5->]\draw[flecha] (c) -- (d);\end{scope}
\end{tikzpicture}
\elabpropia
\end{column}
\begin{column}{0.46\textwidth}\footnotesize
\destacar{En los evangelios}
\begin{itemize}
\item Bautismo: Sal 2:7 + Is 42:1 (Mt 3:17).
\item Is 42:1--4 citado en Mt 12:18--21.
\end{itemize}
\vspace{4pt}
\begin{notalam}[Tradición judía]
El siervo es Israel (Rashi) o el Mesías (Targum de Is 52:13; Chilton, 1987). El Nuevo Testamento reconoce en Jesús ambas cosas.
\end{notalam}
\end{column}
\end{columns}
\end{frame}

\begin{frame}{La vocación en esta etapa}
\centering
\begin{tikzpicture}
 \foreach \x/\n/\s/\t/\d/\r [count=\ovc,evaluate=\ovc as \ovk using int(\ovc+0)] in {0/1/bloque/{Se vivió}/{En el remanente: los Doce, los miles, los sacerdotes. Pablo, «ministro» que ofrece «la ofrenda de las naciones».}/{Ro 15:16; Is 66:20},
   4.3/2/bloqueoro/{Se quebró}/{En la dirigencia del templo y en las facciones de la guerra: la casa de oración convertida en fortaleza.}/{Mt 21:13; \emph{Guerra} 5.1--20},
   8.6/3/bloque/{Se preservó}/{En el remanente y la iglesia de judíos y gentiles; y en Yavne, que custodió la Torá y la oración.}/{Ro 3:2; 11:5}}{\begin{scope}[visible on=<\ovk->]
  \node[\s,anchor=north,text width=3.55cm,minimum height=2.9cm,align=left] (b\n) at (\x,0) {{\small\bfseries\color{ladrillo}\t}\\[4pt]\footnotesize\d\\[4pt]{\scriptsize\color{sepia}(\r)}};\end{scope}}
 \begin{scope}[visible on=<3->]\draw[flecha] (b1.east) -- (b2.west);\end{scope}
 \begin{scope}[visible on=<3->]\draw[flecha] (b2.east) -- (b3.west);\end{scope}
\end{tikzpicture}
\elabpropia
\vspace{8pt}
\begin{clave}
Pablo describe su misión como un oficio sacerdotal (\grc{ἱερουργοῦντα}): «para que la ofrenda de las naciones sea aceptable» (Ro 15:16). La vocación de Éx 19:6 sigue activa en el remanente.
\end{clave}
\end{frame}

\begin{frame}{Los gentiles incorporados}
\begin{columns}[T,onlytextwidth]
\begin{column}{0.48\textwidth}
\centering
\begin{tikzpicture}[font=\scriptsize]
\node[bloqueoro,text width=4.6cm] (a) at (0,0) {\textbf{Antes}: «alejados de la ciudadanía de Israel» (\grc{τῆς πολιτείας τοῦ Ἰσραήλ}, Ef 2:12)};
\begin{scope}[visible on=<2->]\node[bloque,text width=4.6cm] (b) at (0,-1.35) {Cristo derribó «la pared intermedia» (2:14); cf. la balaustrada del templo};\end{scope}
\begin{scope}[visible on=<3->]\node[bloqueazul,text width=4.6cm] (c) at (0,-2.6) {\textbf{Ahora}: «conciudadanos de los santos» (2:19)};\end{scope}
\begin{scope}[visible on=<3->]\draw[flecha] (a) -- (b);\end{scope} \begin{scope}[visible on=<3->]\draw[flecha] (b) -- (c);\end{scope}
\end{tikzpicture}
\elabpropia
\vspace{2pt}
\begin{notalam}[Autoría]
La autoría de Efesios se discute; Romanos sirve de respaldo.
\end{notalam}
\end{column}
\begin{column}{0.48\textwidth}
\begin{cita}Alegraos, naciones, con su pueblo.\fuente{(Ro 15:10; Dt 32:43 LXX)}\end{cita}
\footnotesize Cristo, «servidor de la circuncisión para confirmar las promesas hechas a los padres» (Ro 15:8).\par\smallskip
\begin{clave}\textbf{1~P 2:9}: «real sacerdocio, nación santa» (Éx 19:6). 2:10 cita Os 1:9--10 [2:1], promesa dirigida al reino del norte. Títulos compartidos, no arrebatados a Israel.\end{clave}
\end{column}
\end{columns}
\end{frame}

% ===================================================
\section{Tradición judía}
\begin{frame}{Yohanán ben Zakai y Yavne}
\begin{columns}[T,onlytextwidth]
\begin{column}{0.48\textwidth}\footnotesize
Sacado de Jerusalén en un ataúd, pide a Vespasiano: «Dame Yavne y a sus sabios» (Talmud, Guitín 56a--b; Avot de-Rabí Natán 4).\par
\begin{notalam}[Un relato que creció]
La historia fue creciendo con el tiempo (Neusner, 1970). Su núcleo: en Yavne la vida judía se reorganizó sin templo. La tradición rabínica, escrita entre los siglos \textsc{ii} y \textsc{vi}, no es un relato histórico verificado.
\end{notalam}
\end{column}
\begin{column}{0.48\textwidth}
\begin{cita}Hijo mío, no te aflijas. Tenemos otra expiación como esta: los actos de bondad, como está dicho: Porque misericordia quiero y no sacrificio.\fuente{(Avot de-Rabí Natán A 4; Os 6:6)}\end{cita}
\footnotesize Jesús citó dos veces Os 6:6 (Mt 9:13; 12:7).
\end{column}
\end{columns}
\vspace{4pt}
\centering
\begin{tikzpicture}[font=\scriptsize]
\node[bloqueazul,text width=11.6cm] (m) at (0,0) {El mundo se sostiene sobre tres cosas (Mishná, Avot 1:2)};
\foreach \x/\t in {-3.95/{la Torá},0/{el culto},3.95/{los actos de bondad}}{
  \node[bloque,text width=3.4cm] (c) at (\x,-0.85) {\t};
  \draw[ladrillo,line width=0.6pt] (m.south -| c) -- (c);}
\end{tikzpicture}
\elabpropia
\end{frame}

\begin{frame}{¿Por qué cayó el Segundo Templo?}
\begin{columns}[T,onlytextwidth]
\begin{column}{0.44\textwidth}
\begin{cita}El Segundo Templo, en el que se ocupaban de la Torá, de los mandamientos y de los actos de bondad, ¿por qué fue destruido? Porque había en él odio gratuito.\fuente{(Talmud, Yoma 9b)}\end{cita}
\begin{notalam}[\heb{שִׂנְאַת חִנָּם}, \emph{sinat jinam}]
El Talmud lo ilustra con Kamtsa y Bar Kamtsa (Guitín 55b--56a). Es una autocrítica rabínica, en sintonía con las facciones que describe Josefo.
\end{notalam}
\end{column}
\begin{column}{0.52\textwidth}
\centering
\begin{tikzpicture}[font=\scriptsize]
\node[bloqueazul,text width=2.75cm] (p) at (-1.6,0) {Primer Templo};
\node[bloqueazul,text width=2.75cm] (s) at (1.6,0) {Segundo Templo};
\begin{scope}[visible on=<2->]\node[bloqueoro,text width=2.75cm,minimum height=1.45cm] (pc) at (-1.6,-1.2) {idolatría,\\inmoralidad\\y sangre};\end{scope}
\begin{scope}[visible on=<3->]\node[bloqueoro,text width=2.75cm,minimum height=1.45cm] (sc) at (1.6,-1.2) {Torá y bondad,\\pero odio gratuito};\end{scope}
\begin{scope}[visible on=<3->]\draw[oro!90!ladrillo,line width=1pt,<->] (pc) -- node[above,etiq,text width=0.6cm]{=} (sc);\end{scope}
\begin{scope}[visible on=<4->]\node[bloque,text width=2.75cm] (pd) at (-1.6,-2.45) {destruido (586 a.\,e.\,c.)};\end{scope}
\begin{scope}[visible on=<5->]\node[bloque,text width=2.75cm] (sd) at (1.6,-2.45) {destruido (70 e.\,c.)};\end{scope}
\begin{scope}[visible on=<5->]\draw[flecha] (p) -- (pc);\end{scope} \begin{scope}[visible on=<5->]\draw[flecha] (pc) -- (pd);\end{scope}
\begin{scope}[visible on=<5->]\draw[flecha] (s) -- (sc);\end{scope} \begin{scope}[visible on=<5->]\draw[flecha] (sc) -- (sd);\end{scope}
\begin{scope}[visible on=<6->]\node[etiq,text=sepia,text width=5.8cm] at (0,-3.25) {El odio gratuito «equivale a las tres transgresiones» (Yoma 9b). Ambos caen el 9 de Av (Taanit 4:6).};\end{scope}
\end{tikzpicture}
\notafig{Elaboración propia a partir de Talmud, Yoma 9b.}
\end{column}
\end{columns}
\end{frame}

\begin{frame}{Tisha be-Av: del duelo al consuelo}
\begin{columns}[T,onlytextwidth]
\begin{column}{0.44\textwidth}
\begin{cita}El 9 de Av [\ldots] fue destruida la casa la primera y la segunda vez [\ldots]. Cuando entra Av se disminuye la alegría.\fuente{(Mishná, Taanit 4:6)}\end{cita}
\begin{clave}El consuelo de \emph{Shabat Najamú} es el mismo que esperaba Simeón: «la consolación de Israel» (Lc 2:25; Is 40:1).\end{clave}
\end{column}
\begin{column}{0.52\textwidth}
\centering
\begin{tikzpicture}[font=\scriptsize]
\node[bloqueoro,text width=4.6cm] (a) at (0,0) {\textbf{Entra Av}: «se disminuye la alegría»};
\begin{scope}[visible on=<2->]\node[bloqueoro,text width=4.6cm] (b) at (0,-1.15) {\textbf{9 de Av}: ayuno, Lamentaciones, \emph{kinot}};\end{scope}
\begin{scope}[visible on=<3->]\node[bloque,text width=4.6cm] (c) at (0,-2.3) {El fuego, del 9 al 10 (Taanit 29a)};\end{scope}
\begin{scope}[visible on=<4->]\node[bloqueazul,text width=4.6cm] (d) at (0,-3.45) {\emph{Shabat Najamú}: «Consolad, consolad a mi pueblo» (Is 40:1)};\end{scope}
\begin{scope}[visible on=<4->]\draw[flecha] (a) -- (b);\end{scope} \begin{scope}[visible on=<4->]\draw[flecha] (b) -- (c);\end{scope} \begin{scope}[visible on=<4->]\draw[flecha] (c) -- (d);\end{scope}
\end{tikzpicture}
\elabpropia
\end{column}
\end{columns}
\end{frame}

\begin{frame}{La oración en lugar del sacrificio}
\centering{\large\heb{וּנְשַׁלְּמָה פָרִים שְׂפָתֵינוּ}}\par\smallskip
\begin{cita}Tomad con vosotros palabras y volved a YHWH [\ldots], y te ofreceremos los becerros de nuestros labios.\fuente{(Os 14:2 [14:3])}\end{cita}
\begin{columns}[T,onlytextwidth]
\begin{column}{0.48\textwidth}
\centering
\begin{tikzpicture}[font=\scriptsize]
\node[bloqueazul,text width=3.2cm] (o) at (0,0) {Os 14:2: «becerros de labios»};
\begin{scope}[visible on=<2->]\node[bloque,text width=2.3cm,minimum height=1.25cm] (r) at (-1.5,-1.45) {\textbf{Rabinos}: las oraciones por las ofrendas diarias (Berajot 26b)};\end{scope}
\begin{scope}[visible on=<3->]\node[bloque,text width=2.3cm,minimum height=1.25cm] (h) at (1.5,-1.45) {\textbf{Heb 13:15}, con la LXX: «fruto de labios» (\grc{καρπὸν χειλέων})};\end{scope}
\begin{scope}[visible on=<3->]\draw[flecha] (o.south -| r) -- (r);\end{scope} \begin{scope}[visible on=<3->]\draw[flecha] (o.south -| h) -- (h);\end{scope}
\end{tikzpicture}
\notafig{Elaboración propia. Dos respuestas al mismo verso.}
\end{column}
\begin{column}{0.48\textwidth}
\begin{notalam}[\emph{Birkat ha-minim}: en debate]
Bendición «contra los herejes» compuesta en Yavne (Berajot 28b--29a). ¿Contra los cristianos y con expulsión (Martyn, 1968)? Faltan pruebas para el siglo \textsc{i} (Kimelman, 1981).
\end{notalam}
\end{column}
\end{columns}
\end{frame}

% ===================================================
\section{Restauración}
\begin{frame}{Isaías y Jeremías en esta etapa}
\centering\scriptsize\renewcommand{\arraystretch}{1.18}
\begin{tabular}{>{\RaggedRight\arraybackslash}p{2.4cm}>{\RaggedRight\arraybackslash}p{4.8cm}>{\RaggedRight\arraybackslash}p{4.4cm}}
\toprule
{\color{ladrillo}\bfseries Texto} & {\color{ladrillo}\bfseries Promesa} & {\color{ladrillo}\bfseries Uso en esta etapa}\\\midrule
Is 40:1--3 & Consuelo; camino en el desierto & Juan; Simeón; \emph{Shabat Najamú}\\
Is 49:5--6 & Levantar las tribus; luz de las naciones & Lc 2:32; Hch 1:8; 13:47\\
Is 53:5 & Herido por nuestras rebeliones & Hch 8:32--35; 1~P 2:22--25\\
Is 56:7 & Casa de oración para todos los pueblos & Mc 11:17; Mt 21:13\\
Is 59:20 & El Redentor vendrá a Sion & Ro 11:26 (Estudio 8)\\
Is 61:1--2 & El jubileo del ungido & Lc 4:16--21\\
Jer 7:11 & Juicio sobre el templo & Mt 21:13\\
Jer 31:31--34 [31:30--33] & Nuevo pacto con Israel y Judá & Lc 22:20; Heb 8:8--12\\
Jer 31:35--37 [31:34--36] & Israel no dejará de ser nación & Ro 11:1--2, 29\\
Zac 12:10 & Mirarán al que traspasaron & Jn 19:37; Ap 1:7\\
\bottomrule
\end{tabular}
\elabpropia
\vspace{3pt}
{\footnotesize\itshape\color{sepia}Jeremías prometió el nuevo pacto y la permanencia de Israel mientras caía el Primer Templo: el 70 se entiende a la luz del 586.}
\end{frame}

\begin{frame}{El Redentor vendrá a Sion}
\begin{columns}[T,onlytextwidth]
\begin{column}{0.48\textwidth}
\centering{\large\heb{וּבָא לְצִיּוֹן גּוֹאֵל}}\par\smallskip
\begin{cita}Y vendrá el Redentor a Sion, y a los que se vuelvan de la rebelión en Jacob, dice YHWH.\fuente{(Is 59:20)}\end{cita}
\begin{notalam}[Liturgia judía]
Con este versículo se abre la oración diaria \emph{U-va le-Tsiyón}; la Amidá incorporó la petición diaria por la reconstrucción de Jerusalén (Sacks, 2009).
\end{notalam}
\end{column}
\begin{column}{0.48\textwidth}
\centering
\begin{tikzpicture}[font=\scriptsize]
\node[bloque,text width=4.6cm] (a) at (0,0) {\textbf{Lv 26:44}: «no los desecharé [\ldots] invalidando mi pacto»};
\begin{scope}[visible on=<2->]\node[bloque,text width=4.6cm] (b) at (0,-1.3) {\textbf{Jer 31:36 [31:35]}: la permanencia de Israel, atada al orden del sol y la luna};\end{scope}
\begin{scope}[visible on=<3->]\node[bloqueoro,text width=4.6cm] (c) at (0,-2.6) {\textbf{Is 59:20}: el Redentor viene a la ciudad que el 70 deja en ruinas};\end{scope}
\begin{scope}[visible on=<4->]\node[bloqueazul,text width=4.6cm] (d) at (0,-3.7) {\textbf{Ro 11:26}: Pablo lo citará (Estudio 8)};\end{scope}
\begin{scope}[visible on=<4->]\draw[flecha] (a) -- (b);\end{scope} \begin{scope}[visible on=<4->]\draw[flecha] (b) -- (c);\end{scope} \begin{scope}[visible on=<4->]\draw[flecha] (c) -- (d);\end{scope}
\end{tikzpicture}
\elabpropia
\end{column}
\end{columns}
\end{frame}

\begin{frame}{Daniel 9:26: tres lecturas}
\centering{\large\heb{יִכָּרֵת מָשִׁיחַ וְאֵין לוֹ}}\par\smallskip
\begin{cita}Después de las sesenta y dos semanas se quitará la vida al Mesías, y no tendrá nada; y el pueblo de un príncipe que ha de venir destruirá la ciudad y el santuario.\fuente{(Dn 9:26)}\end{cita}
\vspace{4pt}
\begin{tikzpicture}
 \foreach \x/\n/\t/\d in {0/1/{Cristiana tradicional}/{El ungido es Jesús; Roma destruye ciudad y santuario en el 70.},
   4.3/2/{Histórico-crítica}/{El ungido es Onías~III (c. 171 a.\,e.\,c.); el príncipe, Antíoco~IV.},
   8.6/3/{Judía}/{Otras identificaciones (Rashi). Josefo: Daniel escribió sobre Roma (\emph{Ant.} 10.276).}}{
  \node[caja,anchor=north] (b\n) at (\x,0) {\parbox[t][2.0cm][t]{3.55cm}{\RaggedRight
    {\small\textbf{\textcolor{ladrillo}{\t}}}\\[3pt]\footnotesize \d}};
  \node[circle,fill=cafe,text=piedra,font=\scriptsize\bfseries,inner sep=1.6pt] at (b\n.north west) {\n};}
\end{tikzpicture}
\elabpropia
\end{frame}

\begin{frame}{Perla filológica · los «hasta» de la caída}
\centering\footnotesize\renewcommand{\arraystretch}{1.3}
\begin{tabular}{>{\RaggedRight\arraybackslash}p{1.9cm}>{\RaggedRight\arraybackslash}p{4.0cm}>{\RaggedRight\arraybackslash}p{5.6cm}}
\toprule
{\color{ladrillo}\bfseries Texto} & {\color{ladrillo}\bfseries Griego} & {\color{ladrillo}\bfseries Lo que termina}\\\midrule
Mt 23:39 & \grc{ἕως ἂν εἴπητε} & La ausencia del Mesías para Jerusalén\\
Lc 21:24 & \grc{ἄχρι οὗ πληρωθῶσιν} & La Jerusalén hollada por los gentiles\\
Hch 1:6--7 & \grc{ἐν τῷ χρόνῳ τούτῳ} & La pregunta por el reino: remitida al Padre\\
Hch 3:21 & \grc{ἄχρι χρόνων ἀποκαταστάσεως} & El Mesías en el cielo, hasta la restauración\\
2~Co 3:14--16 & \grc{ἡνίκα ἐὰν ἐπιστρέψῃ} & El velo, cuando Israel se vuelva al Señor\\
Ro 11:25 & \grc{ἄχρι οὗ ... εἰσέλθῃ} & El endurecimiento parcial (Estudio 8)\\
\bottomrule
\end{tabular}
\elabpropia
\vspace{8pt}
\begin{clave}Los mismos textos que anuncian el juicio le fijan un límite: el juicio queda abierto hacia la restauración.\end{clave}
\end{frame}

\begin{frame}{«Tiempos de la restauración»}
\begin{cita}Arrepentíos y volveos [\ldots] para que vengan tiempos de refrigerio, y él envíe al Mesías [\ldots], a quien es necesario que el cielo reciba hasta los tiempos de la restauración de todas las cosas.\fuente{(Hch 3:19--21)}\end{cita}
\vspace{4pt}
\centering
\begin{tikzpicture}[font=\scriptsize]
\foreach \x/\n/\s/\t [count=\ovc,evaluate=\ovc as \ovk using int(\ovc+0)] in {0/1/bloqueoro/{«Arrepentíos y volveos»},3.25/2/bloque/{tiempos de refrigerio},6.5/3/bloque/{envío del Mesías},9.75/4/bloqueazul/{restauración de todas las cosas}}{\begin{scope}[visible on=<\ovk->]
  \node[\s,text width=2.5cm,minimum height=1.0cm] (r\n) at (\x,0) {\t};\end{scope}}
\begin{scope}[visible on=<4->]\foreach \i/\j in {1/2,2/3,3/4}{\draw[flecha] (r\i) -- (r\j);}\end{scope}
\begin{scope}[visible on=<4->]\draw[decorate,decoration={brace,amplitude=4pt,mirror},sepia] (5.2,-0.65) -- (8.0,-0.65) node[midway,below=5pt,etiq]{«el cielo lo recibe \emph{hasta}\ldots»};\end{scope}
\end{tikzpicture}
\elabpropia
\vspace{4pt}
\begin{notalam}[\grc{ἀποκατάστασις}]
Retoma el verbo de la pregunta de Hch 1:6 y apunta a lo que anunciaron los profetas. Pedro llama a sus oyentes «hijos [\ldots] del pacto que Dios hizo con nuestros padres» (Hch 3:25).
\end{notalam}
\end{frame}

\begin{frame}{¿Rechazó Dios a Israel en el 70?}
\begin{columns}[T,onlytextwidth]
\begin{column}{0.44\textwidth}
\begin{cita}Aun estando ellos en tierra de sus enemigos, yo no los desecharé ni los abominaré para consumirlos, invalidando mi pacto con ellos.\fuente{(Lv 26:44)}\end{cita}
\begin{cita}Irrevocables son los dones y el llamamiento de Dios.\fuente{(Ro 11:29)}\end{cita}
\end{column}
\begin{column}{0.52\textwidth}\footnotesize
\begin{enumerate}
\item \textbf{Juicio profético}, como el de Jeremías en el 586, dentro del pacto.
\item \textbf{Los textos tienen un «hasta»}: Mt 23:39; Lc 21:24; Hch 3:21.
\item \textbf{Hubo un remanente}: Pablo, «de la tribu de Benjamín» (Ro 11:1--5).
\item \textbf{El tropiezo tiene propósito}: salvación a los gentiles, celos a Israel (11:11).
\item \textbf{El llamamiento es irrevocable} (11:29), escrito antes del 70 y fundado en Dios.
\end{enumerate}
\begin{clave}Pablo: «tristeza y continuo dolor» (9:2); «mi oración por Israel es para salvación» (10:1). ONE FOR ISRAEL (2025): Israel «ampliado, no reemplazado».\end{clave}
\end{column}
\end{columns}
\end{frame}

% ===================================================
\section{Síntesis}
\begin{frame}{Síntesis}
\begin{tcolorbox}[colback=piedra,colframe=ladrillo,boxrule=1pt,arc=3pt,fontupper=\small]
\textbf{\textcolor{cafe}{Israel de Cristo al año 70:}} el Mesías asume el llamado de su pueblo y establece el nuevo pacto prometido a Israel y Judá; un remanente judío creyente lleva su mensaje «al judío primeramente» y a las naciones, que se suman sin ocupar su lugar. El templo cae en un juicio anunciado, aunque los mismos textos que lo anuncian le fijan un «hasta»; y Pablo, antes del 70, ya lo había dicho: «no ha desechado Dios a su pueblo» (Ro 11:2).
\end{tcolorbox}
\centering\scriptsize\renewcommand{\arraystretch}{1.25}
\begin{tabular}{>{\RaggedRight\arraybackslash}p{3.0cm}>{\RaggedRight\arraybackslash}p{8.6cm}}
\toprule
{\color{ladrillo}\bfseries Eje} & {\color{ladrillo}\bfseries Lo que deja esta etapa}\\\midrule
Mesías de Israel & Gá 4:4; Lc 2:32; Mt 19:28; Lc 22:20; Jn 7; 10:22\\
Vocación & Jesús, el Israel fiel (Mt 2:15; Is 49:3--6); remanente sacerdotal (Ro 15:16)\\
Juicio & Jer 7:11; Mt 23:38; Lc 21:20--24; Josefo; Yoma 9b\\
Restauración & Mt 23:39; Lc 21:24; Hch 1:6; 3:21; Ro 11:1--2, 11, 29\\
\bottomrule
\end{tabular}
\elabpropia
\end{frame}

\begin{frame}{Hacia el Estudio 8}
\begin{columns}[T,onlytextwidth]
\begin{column}{0.46\textwidth}
\begin{cita}El endurecimiento en parte ha acontecido a Israel hasta que haya entrado la plenitud de los gentiles; y así todo Israel será salvo.\fuente{(Ro 11:25--26)}\end{cita}
\footnotesize Si Israel en su mayoría no reconoció al Mesías, ¿qué será de él? Pablo responde en Romanos 9--11.
\end{column}
\begin{column}{0.50\textwidth}
\centering
\begin{tikzpicture}[font=\scriptsize]
\node[bloqueazul,text width=3.6cm,minimum height=1.0cm] (r) at (0,0) {Estudio 8\\Romanos 11: la restauración de Israel};
\foreach \x/\y/\t in {-2.05/1.45/Remanente,0/1.45/Olivo,2.05/1.45/Injerto,-2.05/-1.45/Plenitud,0/-1.45/Misterio,2.05/-1.45/{«Hasta»}}{
  \node[bloqueoro,minimum width=1.7cm] (t) at (\x,\y) {\t};
  \draw[flecha] (t) -- (r);}
\end{tikzpicture}
\elabpropia
\end{column}
\end{columns}
\end{frame}

% ===================================================
\section{Reflexión}
\begin{frame}{Para la reflexión}
\small
\begin{enumerate}
\item Simeón une «luz para las naciones» y «gloria de tu pueblo Israel» (Lc 2:32). ¿Qué se pierde cuando solo se conserva la primera mitad?
\item ¿Por qué importa que Mt 21:43 hable de los dirigentes y de «un pueblo» en singular?
\item Jesús llora sobre Jerusalén antes de anunciar su ruina. ¿Qué actitud pide esto a los creyentes de las naciones (Ro 11:18--21)?
\item Yohanán ben Zakai y Jesús citaron Os 6:6. ¿En qué se parecen y en qué se diferencian sus respuestas a la pérdida del templo?
\item ¿Qué «hasta» de esta etapa le parece más significativo para la esperanza de Israel?
\end{enumerate}
\end{frame}

% ===================================================
\section{Glosario}
\begin{frame}{Glosario}
\begin{columns}[t,onlytextwidth]
\begin{column}{0.49\textwidth}\scriptsize
\textbf{Apokatástasis} (\grc{ἀποκατάστασις}): restauración (Hch 3:21).\par\smallskip
\textbf{Birkat ha-minim} (\heb{בִּרְכַּת הַמִּינִים}): bendición «contra los herejes».\par\smallskip
\textbf{Brit jadashá} (\heb{בְּרִית חֲדָשָׁה}): nuevo pacto (Jer 31:31).\par\smallskip
\textbf{Derór} (\heb{דְּרוֹר}): libertad del jubileo (Is 61:1).\par\smallskip
\textbf{Enkainia} (\grc{τὰ ἐγκαίνια}): Dedicación, Janucá (Jn 10:22).\par\smallskip
\textbf{Gemilut jasadim} (\heb{גְּמִילוּת חֲסָדִים}): actos de bondad.\par\smallskip
\textbf{Kairoí ethnōn} (\grc{καιροὶ ἐθνῶν}): tiempos de los gentiles.
\end{column}
\begin{column}{0.49\textwidth}\scriptsize
\textbf{Kálymma} (\grc{κάλυμμα}): velo (2~Co 3:14--16).\par\smallskip
\textbf{Pléroma} (\grc{πλήρωμα}): plenitud (Gá 4:4; Ro 11:25).\par\smallskip
\textbf{Pórosis} (\grc{πώρωσις}): endurecimiento (Ro 11:25).\par\smallskip
\textbf{Sinat jinam} (\heb{שִׂנְאַת חִנָּם}): odio gratuito (Yoma 9b).\par\smallskip
\textbf{Sukkat David} (\heb{סֻכַּת דָּוִיד}): tienda de David (Am 9:11).\par\smallskip
\textbf{Tisha be-Av} (\heb{תִּשְׁעָה בְּאָב}): 9 de Av, ayuno por los templos.\par\smallskip
\textbf{Yavne} (\heb{יַבְנֶה}): sede de los sabios después del 70.
\end{column}
\end{columns}
\end{frame}

% ===================================================
\section{Referencias}
\begin{frame}[allowframebreaks]{Referencias}
\tiny
\refapa{Allison, D.~C., Jr. (1983). Matt. 23:39 = Luke 13:35b as a conditional prophecy. \emph{Journal for the Study of the New Testament, 5}(18), 75--84.}
\refapa{Avigad, N. (1983). \emph{Discovering Jerusalem}. Thomas Nelson.}
\refapa{Bauckham, R. (1995). James and the Jerusalem church. En R. Bauckham (Ed.), \emph{The book of Acts in its Palestinian setting}. Eerdmans.}
\refapa{Ben-Yehuda, N. (1995). \emph{The Masada myth: Collective memory and mythmaking in Israel}. University of Wisconsin Press.}
\refapa{\emph{Biblia Hebraica Stuttgartensia} (5.ª ed.). (1997). Deutsche Bibelgesellschaft.}
\refapa{Brandon, S.~G.~F. (1951). \emph{The fall of Jerusalem and the Christian church}. SPCK.}
\refapa{Chilton, B.~D. (1987). \emph{The Isaiah Targum} (The Aramaic Bible, Vol. 11). Michael Glazier.}
\refapa{Danby, H. (Trad.). (1933). \emph{The Mishnah}. Oxford University Press.}
\refapa{Davies, W.~D., \& Allison, D.~C. (1997). \emph{A critical and exegetical commentary on the Gospel according to Saint Matthew} (Vol. 3). T\&T Clark.}
\refapa{Epstein, I. (Ed.). (1935--1948). \emph{The Babylonian Talmud}. Soncino Press.}
\refapa{Eusebio de Cesarea. (1926). \emph{The ecclesiastical history} (K. Lake, Trad.; Vol. 1). Harvard University Press.}
\refapa{Evans, C.~A. (2003). \emph{Jesus and the ossuaries}. Baylor University Press.}
\refapa{Fine, S. (2016). \emph{The menorah: From the Bible to modern Israel}. Harvard University Press.}
\refapa{Goldin, J. (Trad.). (1955). \emph{The fathers according to Rabbi Nathan}. Yale University Press.}
\refapa{Goodman, M. (2007). \emph{Rome and Jerusalem: The clash of ancient civilizations}. Allen Lane.}
\refapa{Hays, R.~B. (2016). \emph{Echoes of Scripture in the Gospels}. Baylor University Press.}
\refapa{Jervell, J. (1972). \emph{Luke and the people of God: A new look at Luke-Acts}. Augsburg.}
\refapa{Josefo, F. (1987). \emph{The works of Josephus} (W. Whiston, Trad.). Hendrickson.}
\refapa{Kimelman, R. (1981). Birkat ha-minim and the lack of evidence for an anti-Christian Jewish prayer in late antiquity. En E.~P. Sanders, A.~I. Baumgarten y A. Mendelson (Eds.), \emph{Jewish and Christian self-definition} (Vol. 2). Fortress Press.}
\refapa{Martyn, J.~L. (1968). \emph{History and theology in the Fourth Gospel}. Harper \& Row.}
\refapa{Mason, S. (2016). \emph{A history of the Jewish War, A.D. 66--74}. Cambridge University Press.}
\refapa{Meshorer, Y. (2001). \emph{A treasury of Jewish coins: From the Persian period to Bar Kokhba}. Amphora Books.}
\refapa{Neusner, J. (1970). \emph{Development of a legend: Studies on the traditions concerning Yohanan ben Zakkai}. Brill.}
\refapa{\emph{Novum Testamentum Graece} (28.ª ed.). (2012). Deutsche Bibelgesellschaft.}
\refapa{ONE FOR ISRAEL Ministry. (2025, 3 de julio). \emph{What does «all Israel will be saved» mean in Romans 11?} \url{https://www.oneforisrael.org/what-does-all-israel-will-be-saved-mean-in-romans-11/}}
\refapa{Rahlfs, A., \& Hanhart, R. (Eds.). (2006). \emph{Septuaginta} (ed. rev.). Deutsche Bibelgesellschaft.}
\refapa{Reich, R. (2011). \emph{Excavating the City of David: Where Jerusalem's history began}. Israel Exploration Society.}
\refapa{Sacks, J. (Trad.). (2009). \emph{The Koren Siddur}. Koren Publishers.}
\refapa{Sanders, E.~P. (1985). \emph{Jesus and Judaism}. Fortress Press.}
\refapa{Schürer, E. (1973--1987). \emph{The history of the Jewish people in the age of Jesus Christ} (G. Vermes, F. Millar y M. Goodman, Eds.; ed. rev., 3 vols.). T\&T Clark.}
\refapa{Skarsaune, O., \& Hvalvik, R. (Eds.). (2007). \emph{Jewish believers in Jesus: The early centuries}. Hendrickson.}
\refapa{Snodgrass, K.~R. (2008). \emph{Stories with intent: A comprehensive guide to the parables of Jesus}. Eerdmans.}
\refapa{Tácito. (1931). \emph{Histories, books IV--V; Annals, books I--III} (C.~H. Moore y J. Jackson, Trads.). Harvard University Press.}
\refapa{Tannehill, R.~C. (1985). Israel in Luke-Acts: A tragic story. \emph{Journal of Biblical Literature, 104}(1), 69--85.}
\refapa{Vermes, G. (Trad.). (2004). \emph{The complete Dead Sea Scrolls in English} (ed. rev.). Penguin.}
\refapa{Wright, N.~T. (1996). \emph{Jesus and the victory of God}. Fortress Press.}
\refapa{Yadin, Y. (1966). \emph{Masada: Herod's fortress and the Zealots' last stand}. Random House.}
\end{frame}
\end{document}
